% ============================================================
%  3. JUSTIFICACIÓN   (guía, 2.3)
%  Propósito: explicar POR QUÉ y PARA QUÉ es necesario realizar
%  el proyecto. Debe relacionarse directamente con el problema.
% ============================================================
\chapter{Justificación}
\label{cap:justificacion}

La elaboración del presente trabajo se justifica por la necesidad de contar con una
referencia cuantitativa que indique cuánta demanda de información de transporte debería
registrarse en cada zona del eje metropolitano de Cochabamba según su perfil territorial. Sin
dicha referencia, el registro de consultas de Trufi App, aun siendo extenso, no permite
distinguir las zonas donde la aplicación atiende adecuadamente a su población de aquellas
donde podría existir una necesidad no cubierta.

Desde el punto de vista práctico, una estimación de la demanda esperada permitiría
transformar un registro histórico de uso en un insumo para la toma de decisiones. La
comparación entre lo observado y lo esperado ofrecería a los voluntarios un criterio
adicional, reproducible y verificable para decidir dónde concentrar sus esfuerzos de mapeo,
complementario al conocimiento del territorio que ya poseen. En una organización cuyo recurso
más limitado es el tiempo de sus miembros, incluso una mejora modesta en la asignación de ese
tiempo tiene un impacto apreciable.

Desde la perspectiva de la Ciencia de Datos, el problema presenta un interés técnico
particular. En primer lugar, los datos corresponden a conteos cuya magnitud depende de la
población expuesta, lo que exige técnicas capaces de estimar tasas por habitante y de tratar
la sobredispersión propia de este tipo de variables. En segundo lugar, la demanda de
transporte suele presentar dependencia espacial, de modo que una evaluación convencional
basada en particiones aleatorias podría sobrestimar la capacidad predictiva de cualquier
técnica; por ello, resulta necesario diseñar un esquema de evaluación que respete la
estructura geográfica de los datos. En tercer lugar, la comparación entre técnicas de
distinta complejidad plantea el reto metodológico de seleccionar la solución sin sesgos, de
manera que una mayor sofisticación solo se adopte cuando mejore de forma clara el desempeño.

Los beneficiarios directos del trabajo son Trufi Association y sus voluntarios. Los
beneficiarios indirectos son los usuarios de la aplicación, los gobiernos municipales del eje
metropolitano, los planificadores urbanos y otras comunidades que mantienen planificadores de
rutas en sistemas de paratránsito, para quienes el procedimiento desarrollado podría resultar
replicable.

Es menester destacar que la presente monografía se elabora de manera estructurada y que la
información se desarrolla de forma secuencial, con el propósito de guiar al lector hacia la
resolución del problema planteado. En este sentido, los capítulos iniciales establecen el
problema, el alcance y los objetivos; el marco teórico proporciona los conceptos necesarios;
y el capítulo de desarrollo documenta, fase por fase, cómo los hallazgos de cada etapa
condicionan las decisiones de la siguiente, hasta llegar a los resultados y a su propuesta de
uso.
